% section4_polarization.tex

\section{Observable Implications: Polarization of Ambitious Signals}
\label{sec:polarization}

The binary model delivers a sharp equilibrium statement: once the transition probability into high scrutiny exceeds the critical threshold, the high type exits the ambitious signal and pooling on the safe signal is selected under D1-consistent beliefs. This section records the observable implication in a richer type space. Near the threshold, the quality of the ambitious signal need not decline smoothly. The more distinctive prediction is polarization: ordinary winners withdraw, while low-calibration agents and extreme high types may remain.

\subsection{Continuous types}

Let the type space be an interval
\[
\Theta=[\underline\theta,\overline\theta].
\]
The sender chooses between \(s_A\) and \(s_P\). The current low-scrutiny net separation rent from \(s_A\) is \(B(\theta)\). The true-type-indexed adverse-repricing exposure is \(\Phi(\theta)\), derived from primitive self-accuracy levels
\[
p(\theta)=\Pr(m=1\mid \theta),
\qquad
p'(\theta)>0,
\]
through
\[
\Phi(\theta)=
\alpha D_{\mathrm{KL}}
\left(
\operatorname{Bern}(p(\theta))
\,\middle\|\,
\operatorname{Bern}(\bar p)
\right).
\]
The low-state payoff gain from the ambitious signal is
\[
D(\theta,q)=B(\theta)-\beta q\mu_H\Phi(\theta).
\]
Thus type \(\theta\) chooses \(s_A\) exactly when
\[
q<q^*(\theta),
\qquad
q^*(\theta)=\frac{B(\theta)}{\beta\mu_H\Phi(\theta)}.
\]

The binary theorem corresponds to the case in which the top type's threshold is crossed first. In the continuous model, the entire shape of \(q^*(\theta)\) determines which types remain on the ambitious signal.

\subsection{Hollowing out}

A simple mean-reversion story would predict that, as scrutiny risk rises, the ambitious signal gradually becomes lower quality. That is not the generic implication of top-down pooling. The model removes types according to the exposure-to-rent ratio
\[
\frac{\Phi(\theta)}{B(\theta)},
\]
not according to type itself.

Suppose \(q^*(\theta)\) is nonmonotone: it is low on an intermediate-high region of the type space but rises again in the extreme upper tail. Economically, these are environments in which ordinary winners carry substantial reputation-at-risk but do not have enough surplus to dominate future adverse repricing, while extreme types have sufficiently large current surplus or sufficiently robust self-models to remain exposed.

For a given \(q\), the set of types choosing the ambitious signal is
\[
A(q)=\{\theta\in\Theta:q<q^*(\theta)\}.
\]
If \(q^*(\theta)\) has an interior trough, then for transition probabilities in the corresponding range,
\[
A(q)=[\underline\theta,\theta_1(q)]\cup[\theta_2(q),\overline\theta]
\]
for some \(\theta_1(q)<\theta_2(q)\). The ambitious signal is then generated by two separated groups: low or poorly calibrated types whose exposure-to-rent ratio is low because they have little to lose, and extreme high types whose rents remain large enough to bear exposure.

\begin{proposition}[Polarized degradation]
\label{prop:polarized-degradation}
Suppose \(q^*(\theta)\) is continuous and there exist \(\theta_1<\theta_2\) such that
\[
q^*(\theta)>q \quad\text{for}\quad \theta\in[\underline\theta,\theta_1]\cup[\theta_2,\overline\theta]
\]
and
\[
q^*(\theta)<q \quad\text{for}\quad \theta\in(\theta_1,\theta_2).
\]
Then the type distribution conditional on the ambitious signal has support on two separated intervals. In particular, the ambitious signal is not a monotone high-type signal in this region.
\end{proposition}

\begin{proof}
The sender chooses \(s_A\) if and only if \(q<q^*(\theta)\). The stated inequalities imply that this condition holds exactly on \([\underline\theta,\theta_1]\cup[\theta_2,\overline\theta]\) and fails on \((\theta_1,\theta_2)\). Hence the support of \(\theta\) conditional on observing \(s_A\) is the union of two separated intervals. This support is not an upper set in \(\Theta\). Therefore \(s_A\) is not a monotone high-type signal.
\end{proof}

The observable shadow is a hollowing-out of credible ambition. The conditional mean of type among ambitious senders may fall,
\[
\mathbb E[\theta\mid s_A]\downarrow,
\]
but the more distinctive prediction is a rise in dispersion:
\[
\operatorname{Var}(\theta\mid s_A)\uparrow
\]
relative to the pre-threshold separating region. The surviving ambitious signal can be generated by extreme competence or by poor anticipatory discipline. This is why mean reversion is too weak a description. The ambitious signal becomes polarized.

% Shared figure includes for polarization section
% Keep filenames centralized here for easy updates.

\newcommand{\figPolarizationVarianceLow}{polarization/model-variance-01-v1p21.png}
\newcommand{\figPolarizationVarianceMid}{polarization/model-variance-02-v4p84.png}
\newcommand{\figPolarizationVarianceHigh}{polarization/model-variance-03-v10p24.png}

\newcommand{\figCaptionPolarizationVarianceLow}{\(\operatorname{Var}(\theta\mid s_A)=1.21\)}
\newcommand{\figCaptionPolarizationVarianceMid}{\(\operatorname{Var}(\theta\mid s_A)=4.84\)}
\newcommand{\figCaptionPolarizationVarianceHigh}{\(\operatorname{Var}(\theta\mid s_A)=10.24\)}

\newcommand{\InsertPolarizationVarianceFigure}{%
\begin{figure}[htbp]
\centering
\begin{minipage}{0.32\textwidth}
\centering
\includegraphics[width=\linewidth]{\figPolarizationVarianceLow}
\par\small\figCaptionPolarizationVarianceLow
\end{minipage}
\hfill
\begin{minipage}{0.32\textwidth}
\centering
\includegraphics[width=\linewidth]{\figPolarizationVarianceMid}
\par\small\figCaptionPolarizationVarianceMid
\end{minipage}
\hfill
\begin{minipage}{0.32\textwidth}
\centering
\includegraphics[width=\linewidth]{\figPolarizationVarianceHigh}
\par\small\figCaptionPolarizationVarianceHigh
\end{minipage}
\caption{Model simulations: increasing conditional dispersion for the ambitious signal.}
\label{fig:polarization-variance}
\end{figure}%
}

\InsertPolarizationVarianceFigure

\subsection{Relation to the D1 flip}

The polarization result is the observable analogue of the D1 flip in Section~\ref{sec:dynamic}. In the binary model, once \(q>q_{\mathrm{D1}}\), the off-path ambitious signal is attributed to the low type because the high type has the smaller profitable-deviation response set. In the continuous model, the same force makes the ambitious signal nonmonotone. It no longer means simply ``high type.'' It can mean either extreme surplus or low anticipatory discipline.

The epistemic exposure term remains pinned to true type throughout this argument. The receiver's interpretation affects the reputational return to \(s_A\), while \(\Phi(\theta)\) describes the true-type cost of carrying exposed self-presentation into a possible adverse repricing state. This separation prevents the off-path belief system from redefining the underlying exposure technology.

\subsection{Empirical language}

The model predicts a specific kind of degradation. It is not merely that ambitious signals become noisier. It is that the middle of the credible ambitious class exits first. Agents with too little reputation-at-risk may continue because adverse repricing cannot take much from them. Agents with overwhelming surplus may continue because the ambitious signal remains worthwhile. Agents in between withdraw because they have enough to lose and not enough surplus to dominate the future exposure cost.

In the standard signaling model, the high signal is protected because bad types cannot climb up. In this model, the high signal is hollowed out because good types climb down. The remaining signal can look extreme: true outliers and badly calibrated imitators share the same visible action. That polarization is the observable signature of top-down pooling.

The next step is not to add a richer law of motion. The polarization result already identifies the cross-sectional shadow of withdrawal. What remains is the equilibrium-selection question underneath that shadow: when the ambitious signal has been hollowed out, is high-type withdrawal a self-confirming belief state? Section~\ref{sec:belief-loop} answers this in the binary core. It treats the receiver's interpretation of \(s_A\) and the sender's expectation of that interpretation as one fixed-point object.
